The validation layer compares real-sequence metrics against multiple null and resampling ensembles. Metrics include two-step residual norm, Jensen--Shannon divergence, top singular value, mutual-information excess, low-rank correction improvement, and rank requirements for residual spectral energy. Structure-destroying nulls produce significant separation across several metrics, while block and window bootstrap controls are closer to the observed sequence. This is the desired validation profile: the detected structure is not explained by iid or first-order Markov models, but it is stable under controls that preserve local sequential organization.

The validation should be interpreted conservatively. These tests support a finite computational claim: within the analyzed range and representation, prime residue transitions exhibit higher-order operator structure beyond a first-order Markov baseline. They do not establish an asymptotic theorem about prime gaps or residue transitions.
